\subsection{Two-place witnesses for every regular-prime central matching gate}

\begin{result}{E.212}{Proved}
{A fixed central lifting fibre at a regular odd prime has an exact two-place support test}
Let $\ell$ satisfy the class-number hypothesis of E.209 and let
$w\ge2$. Fix an actual integral class-$(w-1)$ section of the
tripod, with its literal representative retained. Suppose its
global class-$w$ lifting fibre is nonempty, and choose a lift $t$.
At a set $\Sigma$ of finite places $p\ne\ell$, fix actual local
point sections whose class-$(w-1)$ images agree with that
predecessor up to geometric conjugation. Retain their actual
aligning frames as in B.26. Put $b=b_w$ and write
\[
 d_p=d_p(t)\in
 \left(\mathbb Z_\ell/(p^w-1)\mathbb Z_\ell\right)^b
 \tag{E.212.1}
\]
for the central correction required to match the point at $p$.

If $w$ is odd, choose a generator
$\zeta_w\in H^1(G_{\mathbb Q},\mathbb Z_\ell(w))$
and let $\tau_p$ be its local Frobenius value modulo $p^w-1$.
For each $p$ the allowable global correction vectors are either
empty or a coset
\[
 C_p=v_p+\ell^{k_p}\mathbb Z_\ell^b.
 \tag{E.212.2}
\]
Their explicit exponents and representatives are given below.
A simultaneous match at every $p\in\Sigma$ exists if and only
if every $C_p$ is nonempty and
\[
 v_p\equiv v_q\pmod{\ell^{\min(k_p,k_q)}}
       \quad\text{for all }p,q\in\Sigma.
 \tag{E.212.3}
\]
Its failure is therefore witnessed by at most two support places.
If the $k_p$ are unbounded, at most one global correction can
pass all these tests.
If $\Sigma$ is nonempty, they are bounded, and the tests pass, the correction set
is one coset of $\ell^{\max k_p}\mathbb Z_\ell^b$.

If $\ell-1\mid w$, the global correction group is
$(\mathbb Z/\ell^{1+v_\ell(w)})^b$. Every individual local
preimage is empty or a singleton; simultaneous failure again
has a witness at at most two places. If $w$ is even and
$\ell-1\nmid w$, the global correction group is zero and
failure has a one-place witness.

These are tests in one fixed actual lifting fibre. There is no
bound here on the sizes of the witnessing primes or the needed
coefficient exponent, and no assertion that the tests pass or
that their choices can be continued through all weights.
An additional match at the coefficient place is a separate
condition, described below.
\end{result}

\begin{proof}
\emph{The actual correction space and its local coordinates.}
The central kernel of the lift is
$T=\mathbb Z_\ell(w)^b$ in a fixed integral Hall basis.
Inner conjugation acts trivially on this central kernel, and
the arithmetic action is the same scalar $\chi_\ell^w$ on
each of its $b$ coordinates. Consequently global and local
correction maps are the $b$-fold direct sums of the maps for
one Tate module. B.26 identifies the actual global lifting
fibre with the torsor $t+H^1(G_{\mathbb Q},T)$ and gives
the discrepancy classes E.212.1.

There is no additional geometric stabilizer in this integral
calculation. If a nontrivial element fixed by the twisted
predecessor action has first nonzero graded coordinate in
weight $i$, that coordinate is killed by every
$\chi_\ell(g)^i-1$. Globally and at each finite local place
there is a $g$ for which that scalar is nonzero. The graded
coordinate is torsion free, which is a contradiction.
This also proves uniqueness of the actual lower aligning
frame; changing its central lift only changes the resulting
correction by a central coboundary. In particular the
$b$ coordinates cannot be changed independently by extra
lower stabilizers when testing E.212.1.

For odd $w$, E.210 identifies the global parameter space with
$\mathbb Z_\ell^b$. Since every integral weight-$w$ cocycle
is unramified away from $\ell$, localization at $p$ is
\[
 (\lambda_1,\ldots,\lambda_b)\longmapsto
      \tau_p(\lambda_1,\ldots,\lambda_b)
 \quad\text{modulo }p^w-1.
 \tag{E.212.4}
\]
The value $\tau_p$ is independent of the cocycle representative:
a coboundary has value $(p^w-1)c$ at Frobenius. A different
integral generator multiplies every $\tau_p$ by the same unit
and changes the global coordinates inversely. Thus the tests
do not depend on that choice.

\emph{Every odd-weight local preimage is one integral ball.}
Set
\[
 h_p=v_\ell(p^w-1),\qquad
 s_p=\min\{h_p,v_\ell(\tau_p)\},\qquad k_p=h_p-s_p,
 \tag{E.212.5}
\]
where the truncated valuation of the zero class is $h_p$.
Interpret the local module as $R_p^b$, with
$R_p=\mathbb Z/\ell^{h_p}$. The condition
$\tau_p\lambda=d_p$ has a solution precisely when each
coordinate of $d_p$ belongs to $\ell^{s_p}R_p$.
If $s_p=h_p$, this requires $d_p=0$ and leaves every
$\lambda\in\mathbb Z_\ell^b$ allowable; set $k_p=0,v_p=0$.
Otherwise write $\tau_p=\ell^{s_p}u_p$ with
$u_p\in(\mathbb Z/\ell^{k_p})^\times$. Division by
$\ell^{s_p}$ identifies its image with
$\mathbb Z/\ell^{k_p}$, so the exact remaining condition is
\[
 \lambda\equiv
 v_p:=u_p^{-1}(d_p/\ell^{s_p})\pmod{\ell^{k_p}}.
 \tag{E.212.6}
\]
This proves E.212.2, including the individually impossible
and zero-local-group cases.

Two nonempty cosets of this form intersect exactly when
E.212.3 holds. Indeed, if $k_p\le k_q$, intersection is
equivalent to $v_q\equiv v_p$ modulo $\ell^{k_p}$; when
that holds, the whole smaller ball $C_q$ lies in $C_p$.
Thus a pairwise-intersecting family of these cosets is
nested. They are closed subsets of the compact space
$\mathbb Z_\ell^b$, and every finite subfamily has nonempty
intersection. The full intersection is nonempty.
Alternatively, if the exponents are bounded, the largest
one occurs and its coset is contained in all the others;
if they are unbounded, compatible representatives define
one vector by completeness. In the latter case any two
vectors in the intersection differ by every power of
$\ell$, and therefore are equal.

This proves the exact simultaneous criterion and the
two-place witness. The argument uses only the one-scalar
Tate localization in E.212.4; it does not replace an
arbitrary high-rank system of linear equations by a
two-equation test.

\emph{The even weights.}
Suppose $\ell-1\mid w$ and put $a=1+v_\ell(w)$.
By E.210 the generator of global first cohomology is
$\eta_w=(\chi_\ell^w-1)/\ell^a$. Its image at every
$p\ne\ell$ has exact order $\ell^a$.
The $b$-fold localization is therefore injective at each
individual place. The preimage of $d_p$ is empty or a
single vector in $(\mathbb Z/\ell^a)^b$.
If all these preimages are nonempty, they have a common
element exactly when their vectors are all equal. Any
failure is visible either at one empty preimage or at two
different ones. If $w$ is even with $\ell-1\nmid w$,
E.210 gives zero global first cohomology: matching holds
exactly when all $d_p$ are zero, and any failure already
occurs at one place.

\emph{Finite witnesses and their retained data.}
Each local discrepancy and period in these tests belongs
to a finite group. For a witnessing finite set of places
choose $N\ge\max h_p$. Reduction of integral local
$H^1$ to coefficient $\mathbb Z/\ell^N(w)$ is injective,
by the coefficient exact sequence: its kernel is
$\ell^N H^1$, which is zero. The witness is thus visible
at a finite coefficient exponent with the reductions of
the actual lower aligning frames retained. The global
integral correction classes also have finite continuous
reductions. No arbitrary finite-coefficient correction
with additional ramification is inserted into the test.

This proves finite detection of the displayed fixed-fibre
equations. It is not a bound on how many factorial primes
of a changing ordinary pair have to be chosen, nor a
claim that those prime choices by themselves retain the
preceding actual section or these frames. Such a change
requires its own comparison, as in E.202 and E.211.

\emph{The optional coefficient-place match.}
Suppose additionally that an actual point section at
$\ell$ has the prescribed lower image, and use its
unique actual lower aligning frame. In odd weight,
E.210 identifies global and local integral first
cohomology. Its discrepancy consequently selects one
global correction $\lambda_\ell\in\mathbb Z_\ell^b$.
All other matches hold exactly when
$\lambda_\ell\in C_p$ for every $p\in\Sigma$.
Failure is then detected by the coefficient place and
one off-diagonal place. In the nonzero even case,
localization identifies the global correction group
with the local torsion subgroup. The local discrepancy
must lie in that subgroup; if it does, it selects a
unique global correction which must equal all the
support-place corrections. In the remaining even
weights it must be zero. None of these additional
requirements is imposed merely by assuming that the
abelian pair is ordinary at the coefficient place.
\end{proof}

\paragraph{The next gate after E.211.}
For the original minimal-five-shell family and $w=3$,
$b_3=2$, and the local group at $q_j$ or $r_j$ is
$(\mathbb Z/5^{e_j})^2$ with
$e_j=v_5(n_{j+1})$. At two and three the local group
is zero, because neither $2^3-1$ nor $3^3-1$ is
divisible by five. Choose the integral generator
$\zeta_3$ so that its reduction is the epsilon
Kummer class of E.206. Such a generator exists by
E.209 and the coefficient exact sequence.
In the common frames,
\[
                \tau_p\bmod5=\ell_p(\epsilon)
       \quad(p=q_j\text{ or }r_j).
 \tag{E.212.7}
\]
When this residue is nonzero the local preimage is
simply
\[
       \lambda\equiv\tau_p^{-1}d_p\pmod{5^{e_j}},
              \qquad\lambda\in\mathbb Z_5^2.
 \tag{E.212.8}
\]
For the other places E.212.5--6 supplies the exact
division and possible individual obstruction. This
is the explicit next matching test after E.211's
class-two construction, whose initial condition is
realized by E.213, and the global class-three lift
supplied by E.209. The
values $d_p$ are those of that actual lift and
its actual local point comparisons. The existence
of the lift, the regular-prime cohomology theorem,
and finite reciprocity do not assert the pairwise
congruences E.212.3.
